\documentclass[12pt,a4paper]{article}

% --- Paquetes ---
\usepackage[utf8]{inputenc}
\usepackage[T1]{fontenc}
\usepackage[spanish]{babel}
\usepackage{helvet}
\renewcommand{\familydefault}{\sfdefault}
\usepackage[margin=2.5cm]{geometry}
\usepackage{ragged2e}
\usepackage{booktabs}
\usepackage{enumitem}
\usepackage{hyperref}
\usepackage{xcolor}
\usepackage{tabularx}
\usepackage{float}

\justifying

\definecolor{hypercore}{RGB}{41, 98, 255}
\hypersetup{
    colorlinks=true,
    linkcolor=hypercore,
    urlcolor=hypercore
}

\setlength{\parindent}{0pt}
\setlength{\parskip}{0.8em}

\begin{document}

\begin{center}
{\LARGE \textbf{Rifa Digital HyperCore}}\\[0.3em]
{\large ¿Cómo va a funcionar nuestra rifa?}\\[1em]
{\normalsize Equipo HyperCore --- Universidad Tecmilenio}\\
{\normalsize 27 de marzo de 2026}
\end{center}

\vspace{0.5em}
\hrule
\vspace{1em}

% ============================================================
% 1. ¿QUÉ ONDA?
% ============================================================
\section{¿Qué onda con esto?}

Vamos a hacer una \textbf{rifa digital} para juntar lana para el viaje a Cancún. Ganamos la fase regional del reto \textbf{KIA Mexico Innovation MeetUp 2026} en la Universidad Tecmilenio y ahora necesitamos fondear el viaje a la final nacional.

La idea es tener una \textbf{página web} donde la gente pueda:

\begin{itemize}[leftmargin=2em]
    \item Ver de qué se trata la rifa y quiénes somos.
    \item Comprar su boleto digital.
    \item Consultar los resultados cuando se haga el sorteo.
\end{itemize}

Nosotros como equipo vamos a poder:

\begin{itemize}[leftmargin=2em]
    \item Ver quién compró boleto y validar sus pagos.
    \item Ejecutar el sorteo cuando estemos listos.
    \item Notificar a todos los participantes de sus resultados.
\end{itemize}

% ============================================================
% 2. LOS NÚMEROS
% ============================================================
\section{Los números}

\begin{table}[H]
\centering
\begin{tabularx}{\textwidth}{Xr}
\toprule
\textbf{Concepto} & \textbf{Cantidad} \\
\midrule
Costo por persona (avión, hotel, comida, 4 noches) & \$13,000 MXN \\
Personas a fondear & 4 de 5 \\
Total que necesitamos & \$52,000 MXN \\
Lo que queremos juntar con la rifa (la mitad) & \$26,000 MXN \\
Precio del boleto & \$200 MXN \\
Boletos que vende cada quien & 40 \\
Total de boletos & 200 \\
Ingreso bruto si vendemos todo & \$40,000 MXN \\
Fecha límite (aprox.) & 20 de abril de 2026 \\
\bottomrule
\end{tabularx}
\end{table}

La otra mitad (\$26,000) la juntamos por otros medios. Del ingreso de boletos hay que descontar lo que cuesten los premios.

% ============================================================
% 3. PREMIOS
% ============================================================
\section{Los premios}

Habrá \textbf{3 ganadores}:

\begin{table}[H]
\centering
\begin{tabularx}{0.7\textwidth}{cX}
\toprule
\textbf{Lugar} & \textbf{Premio} \\
\midrule
1\textsuperscript{er} lugar & \$5,000 MXN en efectivo \\
2\textsuperscript{do} lugar & Bocina JBL Flip 7 \\
3\textsuperscript{er} lugar & Botella Maestro Dobel \\
\bottomrule
\end{tabularx}
\end{table}

% ============================================================
% 4. ¿CÓMO COMPRA ALGUIEN UN BOLETO?
% ============================================================
\section{¿Cómo compra alguien un boleto?}

Así de sencillo:

\begin{enumerate}[leftmargin=2em]
    \item La persona entra a nuestra página web.
    \item Ve la info de la rifa, los premios, y quiénes somos.
    \item Le da clic a ``Comprar Boleto''.
    \item Llena sus datos: nombre, email, y teléfono (opcional).
    \item Elige si quiere que le mandemos notificaciones por email o no.
    \item Paga por alguno de estos medios:
    \begin{itemize}
        \item \textbf{Efectivo:} nos paga en persona y nosotros confirmamos en el sistema.
        \item \textbf{Transferencia:} sube su comprobante y nosotros lo revisamos.
        \item \textbf{Pago en línea:} paga directo en la página y se confirma solo.
    \end{itemize}
    \item Recibe su boleto digital por email y lo puede ver en la página.
\end{enumerate}

\textbf{Importante:} El boleto no ``cuenta'' para el sorteo hasta que nosotros (o el sistema) confirmemos el pago. Si alguien no paga, su boleto queda pendiente y no entra al sorteo.

% ============================================================
% 5. ¿CÓMO FUNCIONA EL SORTEO?
% ============================================================
\section{¿Cómo funciona el sorteo?}

\begin{enumerate}[leftmargin=2em]
    \item Primero nos aseguramos de que \textbf{todos los pagos estén revisados}. No puede haber boletos pendientes --- todos tienen que estar confirmados o rechazados.
    \item Uno de nosotros (admin) le da al botón de ``Ejecutar Sorteo'' en el panel.
    \item El sistema elige \textbf{3 ganadores al azar} de entre todos los boletos confirmados. Todos tienen la misma probabilidad.
    \item Se publican los resultados en la página.
    \item Se manda un \textbf{email personalizado a cada participante}:
    \begin{itemize}
        \item Si ganó: ``¡Felicidades! Ganaste [premio]. Contáctanos con tu folio para reclamar tu premio.''
        \item Si no ganó: ``Gracias por participar. Te invitamos a seguir nuestro camino en la competencia.''
    \end{itemize}
    \item Los ganadores nos contactan con su número de folio y les entregamos su premio.
\end{enumerate}

No va a ser en vivo. El sistema lo hace automático y publica los resultados.

% ============================================================
% 6. ¿QUÉ PUEDE HACER EL PARTICIPANTE?
% ============================================================
\section{¿Qué puede hacer el participante en la página?}

\begin{itemize}[leftmargin=2em]
    \item \textbf{Ver su boleto} con su folio y datos.
    \item \textbf{Actualizar su info} (nombre, email, teléfono).
    \item \textbf{Consultar resultados} después del sorteo (por folio o email).
    \item \textbf{Salirse de la rifa} si ya no quiere participar. No hay reembolso automático, pero si nos contacta podemos ver cómo ayudarle.
    \item \textbf{Pedir que borren sus datos} si quiere que eliminemos su información.
\end{itemize}

% ============================================================
% 7. ¿QUÉ HACEMOS NOSOTROS COMO EQUIPO?
% ============================================================
\section{¿Qué hacemos nosotros como equipo?}

Desde un panel de administración vamos a poder:

\begin{itemize}[leftmargin=2em]
    \item Ver todos los boletos vendidos y su estado (pendiente, confirmado, rechazado).
    \item \textbf{Revisar comprobantes} de transferencia y aceptar o rechazar pagos.
    \item \textbf{Confirmar pagos en efectivo} manualmente.
    \item \textbf{Ejecutar el sorteo} cuando todo esté listo.
    \item Ver los resultados y el estado de las notificaciones enviadas.
\end{itemize}

% ============================================================
% 8. LA PÁGINA "NOSOTROS"
% ============================================================
\section{La página ``Nosotros''}

Esta sección es súper importante para que la gente confíe en la rifa. Va a incluir:

\begin{itemize}[leftmargin=2em]
    \item Quiénes somos --- presentación del equipo HyperCore.
    \item Foto y rol de cada integrante.
    \item Links a nuestros perfiles de LinkedIn.
    \item Links a los certificados del Innovation MeetUp (para que puedan verificar).
    \item Una breve explicación del reto de KIA y por qué necesitamos los fondos.
    \item Links a nuestras redes sociales.
    \item Email de contacto.
\end{itemize}

La idea es que cualquier persona que entre y no nos conozca pueda verificar que somos reales y que la causa es legítima.

% ============================================================
% 9. FORMAS DE PAGO
% ============================================================
\section{Formas de pago}

\begin{table}[H]
\centering
\begin{tabularx}{\textwidth}{lXl}
\toprule
\textbf{Método} & \textbf{¿Cómo funciona?} & \textbf{¿Quién confirma?} \\
\midrule
Efectivo & Nos pagan en persona & Nosotros en el panel \\
Transferencia & Suben comprobante en la página & Nosotros revisamos \\
Pago en línea & Pagan directo en la página & Se confirma solo \\
\bottomrule
\end{tabularx}
\end{table}

La pasarela de pago en línea todavía está por definirse. Hay que revisar comisiones de Mercado Pago, Stripe, PayPal, etc.

% ============================================================
% 10. NOTIFICACIONES
% ============================================================
\section{Notificaciones por email}

El email es \textbf{opcional}. El participante decide si quiere recibirlos o no.

\begin{itemize}[leftmargin=2em]
    \item \textbf{Confirmación de boleto:} cuando validamos su pago.
    \item \textbf{Resultado del sorteo:} personalizado --- si ganó o no, con invitación a seguirnos.
\end{itemize}

Si alguien no quiere emails, puede consultar resultados en la página o en nuestras redes. Y si gana, lo contactamos personalmente.

% ============================================================
% 11. ¿QUÉ PASA SI ALGO SALE MAL?
% ============================================================
\section{¿Qué pasa si algo sale mal?}

Ya pensamos en varios escenarios:

\begin{itemize}[leftmargin=2em]
    \item \textbf{Alguien no paga:} su boleto queda pendiente y no entra al sorteo. Si pasa mucho tiempo, se cancela automáticamente.
    \item \textbf{No llega la confirmación de un pago en línea:} el sistema revisa periódicamente con la pasarela para verificar.
    \item \textbf{Intentamos hacer el sorteo sin boletos confirmados:} el sistema no lo permite.
    \item \textbf{Falla el envío de un email:} se reintenta automáticamente varias veces. Si sigue fallando, nos avisa para revisarlo manualmente.
    \item \textbf{Alguien intenta ejecutar el sorteo dos veces:} el sistema no lo permite. Los resultados no se pueden cambiar una vez publicados.
\end{itemize}

% ============================================================
% 12. PENDIENTES
% ============================================================
\section{Lo que falta por definir}

Estos son los puntos que necesitamos resolver como equipo:

\begin{table}[H]
\centering
\begin{tabularx}{\textwidth}{cX}
\toprule
\textbf{\#} & \textbf{Pendiente} \\
\midrule
1 & Revisar la ley de protección de datos (LFPDPPP) --- ¿qué necesitamos cumplir? \\
2 & Elegir la plataforma de pago en línea (revisar comisiones) \\
3 & Confirmar la fecha exacta del sorteo (20 de abril es aproximado) \\
4 & Definir las herramientas para construir la página \\
5 & Conseguir los links de los certificados del Innovation MeetUp \\
6 & Definir qué redes sociales vamos a vincular \\
7 & Elegir el servicio para enviar emails \\
8 & Decidir dónde vamos a hospedar la página \\
\bottomrule
\end{tabularx}
\end{table}

% ============================================================
% CIERRE
% ============================================================
\vspace{1.5em}
\hrule
\vspace{1em}

\begin{center}
\textit{Este documento resume la idea de cómo va a funcionar nuestra rifa digital. Conforme vayamos resolviendo los pendientes, lo actualizamos. Si tienen dudas, ideas o sugerencias, díganle al equipo.}\\[1em]
{\Large \textbf{¡Vamos a Cancún!}} $\star$
\end{center}

\end{document}
